%! Factoring: GCF and Grouping — Unit 1, Lesson 1.3 — Homework (Answer Key)
\documentclass[10pt]{article}
\usepackage{atda-article}
\usepackage{atda-key}   % pulls in -boxes; do NOT also load -boxes
\newcommand{\TallMath}[1]{$\displaystyle #1\rule[-1.4em]{0pt}{3.2em}$}
\setlength{\workrowsep}{8pt}

\begin{document}

\pageheader{Unit 1, Lesson 1.3}{Homework --- Answer Key}
% NAMESTRIP: no \namedateperiod here — mirrors the blank. Teacher-only prose goes
% in the lesson plan as \begin{teachernote}[Homework], never in a key.

\vspace{0.15in}

\boxguard[16]
\begin{notesbox}{Practice --- GCF and Grouping}
\small
Show your work. Factor \textbf{completely} every time: when you are done, the terms inside each
factor should share nothing but $1$. Check by multiplying back.

\begin{enumerate}[leftmargin=1.6em, itemsep=6pt, topsep=4pt]

\item Factor: $12x+30$
\begin{work}
  12x+30 &= 6(2x) + 6(5) \\
         &= 6(2x+5)
\end{work}

\item Factor: $9x^3-15x^2$
\begin{work}
  9x^3-15x^2 &= 3x^2(3x) - 3x^2(5) \\
             &= 3x^2(3x-5)
\end{work}

\item Factor: $8a^2b^2+20ab^3$
\begin{work}
  8a^2b^2+20ab^3 &= 4ab^2(2a) + 4ab^2(5b) \\
                 &= 4ab^2(2a+5b)
\end{work}

\item Factor: $16y^4-24y^3+40y^2$
\begin{work}
  16y^4-24y^3+40y^2 &= 8y^2\left(2y^2\right) - 8y^2(3y) + 8y^2(5) \\
                    &= 8y^2\left(2y^2-3y+5\right)
\end{work}

\item Factor out a \emph{negative} GCF: $-6x^2-15x$
\begin{work}
  -6x^2-15x &= -3x(2x) - 3x(5) \\
            &= -3x(2x+5)
\end{work}

\item Factor by grouping: $x^3+6x^2+4x+24$
\begin{work}
  x^3+6x^2+4x+24 &= x^2(x+6) + 4(x+6) \\
                 &= (x+6)\left(x^2+4\right)
\end{work}

\item Factor by grouping: $10x^3-15x^2+8x-12$
\begin{work}
  10x^3-15x^2+8x-12 &= 5x^2(2x-3) + 4(2x-3) \\
                    &= (2x-3)\left(5x^2+4\right)
\end{work}

\item Factor by grouping (watch the sign): $x^3-2x^2-7x+14$
\begin{work}
  x^3-2x^2-7x+14 &= x^2(x-2) - 7(x-2) \\
                 &= (x-2)\left(x^2-7\right)
\end{work}

\item Factor completely --- GCF first, then group: $3x^3+6x^2+15x+30$
\begin{work}
  3x^3+6x^2+15x+30 &= 3\left(x^3+2x^2+5x+10\right) \\
                   &= 3\left[x^2(x+2) + 5(x+2)\right] \\
                   &= 3(x+2)\left(x^2+5\right)
\end{work}

\end{enumerate}
\end{notesbox}

\vspace{0.10in}

\boxguard[16]
\begin{scenariobox}[In Context --- The Shipping Crate]{royal}
\small
A packing company builds a rectangular crate whose volume is
\[ V(x) = 8x^3+12x^2+10x+15 \quad \text{cubic feet,} \]
where $x$ is the size setting on the cutting machine, in feet.

\begin{enumerate}[leftmargin=1.6em, itemsep=6pt, topsep=4pt, start=10]

\item Factor $V(x)$ completely by grouping.
\begin{work}
  8x^3+12x^2+10x+15 &= 4x^2(2x+3) + 5(2x+3) \\
                    &= (2x+3)\left(4x^2+5\right)
\end{work}

\item The crate's height is $(2x+3)$ feet. Write the expression for the area of its base, then
find the height, the base area, and the volume when $x=3$. Include units.
\begin{work}
  \text{base area} &= 4x^2+5 \\
  x=3: \quad 2x+3 &= 9 \text{ ft}, \qquad 4x^2+5 = 41 \text{ sq ft} \\
  V &= 9 \cdot 41 = 369 \text{ cubic feet}
\end{work}

\item In one sentence, say what the factored form tells the packing company that the expanded form
does not.

\ansline{It gives the crate's actual dimensions --- height $(2x+3)$ ft and base area $\left(4x^2+5\right)$ sq ft.}

\end{enumerate}
\end{scenariobox}

\vspace{0.10in}

\boxguard[16]
\begin{scenariobox}[A First Look Ahead]{goldacc}
\small
\begin{enumerate}[leftmargin=1.6em, itemsep=6pt, topsep=4pt, start=13]

\item Grouping needs four terms --- but a three-term polynomial can be \emph{turned into} four.
Factor $x^2+7x+10$ this way: find two numbers that multiply to $10$ and add to $7$, use them to
split the middle term $7x$ into two terms, then group. This is exactly how Lesson 1.4 factors
trinomials.
\begin{work}
  x^2+7x+10 &= x^2+5x+2x+10 \\
            &= x(x+5) + 2(x+5) \\
            &= (x+5)(x+2)
\end{work}

\end{enumerate}
\end{scenariobox}

\vspace{0.10in}

\boxguard[18]
\begin{extensionbox}
\small
\textbf{When the common factor is not a monomial.} A rectangle's area is written
$A(x) = x(x+4)+3(x+4)$. Both pieces share the whole binomial $(x+4)$, so it can be pulled out front
just like a monomial GCF. Factor it, then multiply your answer back out. Why does treating a
binomial as ``the shared piece'' work for exactly the same reason a monomial does?
\begin{work}
  x(x+4)+3(x+4) &= (x+4)(x+3) \\
                &= x^2+7x+12 \quad\checkmark
\end{work}
\ansline{The distributive property never asked what the shared factor was --- only that it be shared.}
\end{extensionbox}

\vspace{0.10in}

\begin{remindbox}
\small
\textbf{Next: Lesson 1.4 --- Factoring Trinomials} (\texttt{A2.EO.3b}). Item 13 already ran the
engine: split the middle term, then group. Lesson 1.4 gives you a reliable way to find the two
numbers to split it with, so three-term polynomials stop being guesswork.
\end{remindbox}

\end{document}
